\section{A Model of Liquid Democracy}\label{sec:model}

In this section, we present the liquid democracy game. This game generalizes the game-theoretic 
model studied by Escoffier et al.~\cite{EGP19}. 

\subsection{The Delegation Graph}\label{sec:delegation-graph}
In liquid democracy each agent has three strategies: she can abstain, vote herself, or delegate her 
vote to another agent. So we can represent an instance by a directed network $G=(V,A)$
called the {\em delegation graph}.
There is a vertex in $V=\{1,\cdots, n\}$ for each agent. 
To define the sets of arcs, there are three possibilities.
First, if agent $i$ votes herself the delegation graph contains a self-loop $(i,i)$.
Second, if agent $i$ delegates her vote to agent $j\neq i$ 
then there is an arc $(i,j)$ in $G$.
Third, if agent $i$ abstains then the vertex $i$ has out-degree zero.

Now, because the out-degree of each vertex is at most one, the delegation graph $G$ is a $1$-forest.  
That is, each component of $G$ is an arborescence plus at most one arc.
In particular, each component is either an arborescence and, thus, contains no cycle or 
contains exactly one directed cycle (called a {\em delegation cycle}).
In the former case, the component contains one sink node corresponding to an abstaining voter.
In the latter case, if the delegation cycle is a self-loop the component contains exactly one voter called
a {\em guru}; if the delegation cycle has length at least two then the component contains no voters. 
An example of a delegation graph is shown in Figure~\ref{fig:delegation-graph}.
\begin{figure}
	\begin{center}
	\includegraphics*[scale=0.13]{ExampleLD}
	\end{center}	
	\caption{A delegation graph.}\label{fig:delegation-graph}	
\end{figure}

Observe that, by the transitivity of delegations, if an agent $i$ is in a component containing
a guru $g$ then that guru will cast a vote on $i$'s behalf. On the other hand, if agent $i$ is 
in a component without a guru (that is, with either a sink node or a cycle of length at least two)
then no vote will be cast on $i$'s behalf. 
We denote the guru $j$ representing agent $i$ by $g(i)=j$ if it exists (we write $g(i)=\emptyset$ otherwise).
Furthermore, its easy to find $g(i)$: simply apply path traversal starting at vertex $i$ in
the delegation graph. 

For example, in Figure~\ref{fig:delegation-graph} two components contain a guru. 
Agent $4$ is the guru of agents $1,2,3$ and itself; agent $9$ is the guru only for itself. 
The vertices in the remaining three components have no gurus.
There are two components with {\em delegation cycles}, namely $\{(10,11), (11, 10)\}$ and 
$\{(5,6), (6, 7), (7,5)\}$. The final component also contains no guru as agent $12$ is a sink node and thus abstains.






\subsection{The Liquid Democracy Game}\label{sec:delegation-game}
The game theoretic model we study is a generalization of the model of Escoffier et al.~\cite{EGP19}.
A pure strategy $s_i$ for agent $i$ corresponds to the selection of at most one outgoing arc.
Thus we can view $s_i$ as an $n$-dimensional vector $\x_i$. Specifically, 
$\x_i$ is a single-entry vector with entry $x_{ij}=1$ if $i$ delegates to agent $j$ and $x_{ij}=0$ otherwise.
Note that if $x_{ii}=1$ then $i$ votes herself (``delegates herself'') and that $\x_i={\bf 0}$ if agent $i$ abstains.


It immediately follows that there is a unique delegation graph $G_{\x}$ associated with
a pure strategy profile $\x=(\x_1, \x_2 \cdots, \x_n)$.
To complete the description of the game, we must define the payoffs corresponding to
each pure strategy profile $\x$. To do this, let agent $i$ have a utility $u_{ij}\in [0,1]$ if she 
has agent $j$ as her guru. Because of the costs of voting in terms of time commitment, knowledge acquisition, etc.,
it may be that $u_{ii}< u_{ij}$ for some $j\neq i$ trusted by $i$.\footnote{Indeed, if this is not the case then
liquid democracy has no relevance.}
We denote the utility of agent $i$ in the delegation graph $G_{\x}$ by $u_i(\x)=u_{i,g(i)}$. 
If agent~$i$ has no guru then it receives zero utility.\footnote{We remark that our results hold even when agents who abstain or have no guru obtain positive utility.}
It follows that only agents that lie in a component of $G_{\x}$ containing a guru can obtain positive utility.

For example, in Figure~\ref{fig:delegation-graph} agent $9$ is a guru so receives utility $u_{9,9}$.
Each agent $i\in \{1,2,3,4\}$ has agent $4$ as its guru so receives utility $u_{i,4}$. 
All the remaining agents have no guru and so receive zero utility.   

An agent $i$ is playing a best response at a pure strategy $\x=(\x_1, \x_2, \cdots, \x_n)$ if he cannot increase
his utility by selecting a different or no out-going arc. The strategy profile is a pure Nash equilibrium if every agent is
playing a best response at $\x$. 
Our interest is in comparing the social welfare of equilibria to the optimal welfare in the game.
To do this, let the social welfare of $\x$ be $\SW(\x) = \sum_{i\in V} u_i(\x)$ and let 
$\Opt=\max\limits_{\x} \sum_{i\in V} u_i(\x)$ be the optimal welfare over all strategy profiles. 
The {\em price of stability} is the worst ratio over all instances between the {\em best} welfare of a Nash equilibrium
and the optimal social welfare. %{\bf [Formal definition?...]}

The reader may ask why could an equilibrium have low social welfare. The problem is that delegation is transitive, but {\em trust is not}.
 Agent $i$ may delegate to an agent $j$ where $u_{ij}$ is large but $j$ may then re-delegate to an agent $k$ where $u_{ik}$ is small.
Worse, agent $i$ may receive no utility if the transitive delegation of its vote leads to a delegation cycle or an abstaining voter.
 Unfortunately, not only can pure Nash equilibria have low social welfare in the liquid democracy game they need not even exist!
 \begin{lem}\label{lem:no-PSNE}
There exist liquid democracy games with no pure strategy Nash equilibrium. 
\end{lem}
\begin{proof}
Let there be three voters with $u_1=(\frac{1}{2},1,0)$, $u_2=(0,\frac{1}{2},1)$ and $u_3=(1,0,\frac{1}{2})$. 
(A similar instance was studied in \cite{EGP19}.) Assume $\x$ is a pure Nash equilibrium and let $S$ be the set of gurus 
in $G_{\x}$. There are two cases.
First, if $|S|\leq 1$ then there exists an agent $i$ with zero utility. This agent can deviate and vote herself to 
obtain utility of $\frac{1}{2}>0$, a contradiction. 
Second, if $|S|\geq 2$ then one of the gurus can delegate its vote to another guru to obtain a utility of $1> \frac{1}{2}$, a contradiction.
Therefore, no pure Nash equilibrium exists. \qed 
\end{proof}


\subsection{Mixed Strategy Equilibria}\label{sec:MSNE}

Lemma~\ref{lem:no-PSNE} tells us that to obtain performance guarantees we must look beyond pure strategies.
Of course, as liquid democracy is a finite game, a mixed strategy Nash equilibrium always exists.
It follows that we can always find a mixed Nash equilibrium in the liquid democracy game and compare
its welfare to the optimal welfare.

Here a mixed strategy for agent $i$ is now simply a non-negative vector $\x_i$ where the entries 
satisfy $\sum_{j=1}^n x_{ij}\le 1$. Note that agent $i$ then abstains with probability $1-\sum_{j=1}^n x_{ij}$.
However, because abstaining is a weakly dominated strategy, we may assume no voter abstains in either the
optimal solution or the best Nash equilibria. Hence, subsequently, each $\x_i$ will be unit-sum vector.

What is the utility $u_i(\x)$ of an agent for a mixed strategy profile $\x=(\x_1, \x_2, \dots, \x_n)$?
It is simply the expected utility  given the probability distribution over the delegation graphs generated by $\x$.
%Namely
%$$u_i(\x)= \sum_{G} \mathbb{P}_\x(G)\cdot u_i(G)$$
%Here $\mathbb{P}_\x(G)=\prod_{(i,j)\in G} x_{ij}$ because $\x_i$ is unit-vector and so every vertex has exactly one
%outgoing arc in each delegation graph generated. Further $u_i(G)=u_{i, g^G(i)}$ where $g^G(i)$ is the guru of $i$ in the
%graph $G$.
There are several equivalent ways to define this expected utility. For the purpose of analysis, 
the most useful formulation is in terms of directed paths in the delegation graphs.
Denote by $\mathcal{P}(i,j)$ be the set of all directed paths from $i$ to $j$ in a complete directed graph 
on $V$. Let $\mathbb{P}_{\x}( g(i)=j)$ be the probability that $j$ is the guru of $i$ given the mixed strategy profile $\x$. Then
\begin{align*}
	u_i(\x) &= \sum_{j\in V} \mathbb{P}_{\x}( g(i)=j  ) \cdot u_{ij} \\ 
	 &= \sum_{j\in V} \mathbb{P}_{\x}( \exists \text{ path from } i \ \text{to} \ j  )\cdot \mathbb{P}_{\x}(\exists \text{ arc } (j,j) ) \cdot u_{ij}\\
	&=\sum_{j\in V} \left( \sum_{P \in \mathcal{P}(i,j)}  \prod\limits_{a\in P}  x_a \right)  \cdot x_{jj} \cdot u_{ij}
\end{align*}
To understand this recall that $\x_i$ is a unit-sum vector. It follows that each delegation graph generated by
the mixed strategy $\x$ has uniform out-degree equal to one.
This is because we generate a delegation graph from $\x$ by
selecting exactly one arc emanating from $i$ according to the probability distribution $\x_i$.
Consequently, $j$ can be the guru of $i$ if and only if the delegation graph contains a self-loop $(j,j)$ {\em and} contains 
a unique path $P$ from $i$ to $j$. The second equality above then holds.
Further, the choice of outgoing arc is independent at each vertex $i$. This implies the third equality.

Regrettably, however, no welfare guarantee is obtainable even with mixed strategy Nash equilibria.

 \begin{lem}\label{lem:bad-MSNE}
The price of stability in liquid democracy games is zero. 
\end{lem}
\begin{proof}
	
	Consider an instance with three agents whose utility is given as $u_1=(\delta,1,0),u_2=(0,\delta,1) $ and $u_3=(1,0,\delta)$. The optimal welfare is $\Opt=1+2\delta$ obtained by the first and second agent voting and the third agent delegates to the first agent.  
	
	By a similar argument to Lemma~\ref{lem:no-PSNE}, we know that no pure strategy Nash equilibrium exists. It is straightforward to verify that this game has a unique mixed strategy Nash equilibria $\x$, where $x_1=(\delta,1-\delta,0)$, $x_2=(0,\delta,1-\delta)$ and $x_3 =(1-\delta, 0, \delta) $. 
	
 The expected utility of agent~$1$ is then
\begin{eqnarray*}
u_1(\x) &=& \sum_{j\in V} \left( \sum_{P \in \mathcal{P}(i,j)}  \prod\limits_{a\in P}  x_a \right)  \cdot x_{jj} \cdot u_{ij} \\
&=& x_{11} \cdot u_{11} + x_{12}\cdot x_{22} \cdot u_{12}  \\
&=& \delta^2 + (1-\delta)\cdot \delta \\
&=& \delta
 \end{eqnarray*}
The case of agents $2$ and $3$ are symmetric, hence the social welfare of this equilibria is $3\delta$. 
Thus, the price of stability is $\frac{3\delta}{1+2\delta}$ which tends to zero
as $\delta\rightarrow 0$. \qed 
\end{proof}

Lemma~\ref{lem:bad-MSNE} appears to imply that no reasonable social welfare guarantees
can be obtained for liquid democracy. This is not the case. Strong performance guarantees
can be achieved, provided we relax the incentive constraints.
Specifically, we switch our attention to approximate Nash equilibria.
A strategy profile $\x$ is an $\epsilon$-{\em Nash equilibrium} if, for each agent $i$, 
$$u_{i}(\x)\geq \ (1-\epsilon)\cdot u_{i}(\hat{\x}_i,\x_{-i}) \qquad \forall \hat{\x}_i $$

Above, we use the notation $\x_{-i}=\{\x_i\}_{j\neq i}$.
Can we obtain good welfare guarantees for approximate Nash equilibria? We will prove the answer is {\sc yes}
in the remainder of the paper. In particular, we present tight bounds on the price of stability for $\epsilon$-Nash equilibria.
%\textbf{Furthermore we show the bounds obtained therein are tight."Should we include that our result is tight here"?}

